\section{Research Experience}

\resumeExpEntry
    {Graduate Researcher}
    {Aug 2026 – Present}
    {Department of Physics \& Astronomy, University of Delaware}
    {Newark, DE}
    {Advisor: Prof. Sarah Dodson-Robinson}
    {
        \item Analyzing solar observations to test and calibrate spectral coherence methods for measuring stellar rotation and magnetic activity cycles.
        \item Developing dual-frequency coherence methods to disentangle stellar activity from planetary signals in radial velocity (RV) datasets.
    }

\resumeExpEntry
    {Undergraduate Visiting Researcher}
    {Jun 2025 – Jul 2026}
    {Department of Physics \& Astronomy, University of Delaware}
    {Newark, DE}
    {Advisor: Prof. Sarah Dodson-Robinson}
    {
        \item Evaluated advanced statistical methods to distinguish genuine exoplanetary signals from stellar noise in high-precision RV data.
        \item Implemented gaussian process (GP) regression and frequency-domain analysis to mitigate stellar variability and improve the reliability of RV detections.
        \item Analyzed 25 years of multi-instrument RV datasets to determine the planetary architecture of Barnard's Star, utilizing advanced statistical metrics for model validation.
    }

\resumeExpEntry
    {Undergraduate Researcher}
    {Nov 2024 – May 2025}
    {Department of Physics, Universidad del Valle de Guatemala}
    {Guatemala City}
    {Advisor: Prof. Eduardo Rubio-Herrera}
    {
        \item Processed and analyzed photometric time series data from the TESS mission to identify exoplanet candidates.
        \item Developed an algorithm to detect transits and extract orbital parameters from stellar light curves.
    }
